\section{Introduction}\label{sec:vtm-introduction}

Let $\mathbb P_{n;q,\boldsymbol\mu}^{\rm af}$ be the all-flux
chamber QED law on $C(n)=\Z^3/n\Z^3$, with charge $q>0$,
real lattice mass $\boldsymbol\mu$, Wilson parameter $0<\rho<1$,
weights $(6,6\sqrt2,6\sqrt2,6)$ and antiperiodic fermions.
Its real curvature is $F^{\rm u}=dA+(2\pi/q)z$, where $z$ is
a closed integer two-cochain. Write $Z_n$ for its unnormalized
partition function and $F^\sharp$ for the physical step field:
on the cube of side $\eta$ centered at a plaquette, its
corresponding component is $\eta^{-3/2}F^{\rm u}$.
The measure and its normalization are recalled in
\cref{sec:vtm-data}.

Fix an odd integer $L\ge13$, $M=L^k$ with integer $k\ge2$,
$0<\varepsilon\le1/200$ and $0<\bar m\le1$. Set
\begin{equation}\label{eq:cutoff-scales}
 h=L^{\lfloor\log_L(1/\bar m)\rfloor},\qquad
 \eta=L^{-N},\quad q=e\sqrt\eta,\qquad R=8hL^T,\quad n=R/\eta.
\end{equation}
At each $e,N$ let $\boldsymbol\mu_N$ be real and independent
of the outer volume $T$. Every reference and outer law below
uses this same mass. The cutoffs $N,T$ are nonnegative integers.
For $e\sqrt h<1$, put
\[
 v_K=-\log(e\sqrt h),\qquad
 r_K=\min\{L^i:i\in\Z_{\ge0},\ L^i\ge v_K^2\},
 \qquad R^*=\frac{hMr_K}{L}=hL^{T^*}.
\]
The reference torus has side $n_*=R^*/\eta$ and
$N_*=8L^3$ cells, each of physical width $\ell_c=R^*/(2L)$.

Fix real nonzero $\beta\in C_c^\infty(\R^3;\Lambda^2)$ and
$f=dd^*\beta$. Let $D_\beta$ be the side of a least coordinate
cube containing $\supp\beta$ and let $Q_\beta$, of side
$D_\beta^{\rm buf}=D_\beta+4$ and center $z_\beta$, be that
cube enlarged by two in each direction. Denote cube averaging
by $J_\eta$ and centered differences by
\[
 \delta_{\eta,\mu}g=
 \frac{g(\,\cdot+\eta e_\mu/2)-g(\,\cdot-\eta e_\mu/2)}{\eta}.
\]
Define
\begin{equation}\label{eq:cells-smooth-test}
 a'_{\eta,\nu}=-\sum_{\mu\ne\nu}\delta_{\eta,\mu}\beta_{\mu\nu},
 \qquad
 f^\eta_{\mu\nu}=\delta_{\eta,\mu}a'_{\eta,\nu}
                     -\delta_{\eta,\nu}a'_{\eta,\mu}.
\end{equation}
The covering bond source is
$\widehat a^o_\nu(x)=\eta^{1/2}(J_\eta a'_{\eta,\nu})(\eta x+\eta e_\nu/2)$.
Put $v_\eta=\|d\widehat a^o\|_{\ell^2}^2$.
A centered placement $v\in\eta\Z^3$ satisfies
$|v+z_\beta-(\ell_c/2)(1,1,1)|_\infty\le\eta/2$.
Translate the probe by $v$, reflect its source into the
$N_*$ cells, and periodize to obtain $\widehat a_c$.
The reflected copies use the two-form orientation signs.

Let $D^{\rm ch}[B]$ be the chamber Dirac operator with bond
phase $B$. Define the free one-cell ratio and its ordinary
derivatives by
\begin{equation}\label{eq:response-coefficient-relations}
 N_0(t)=\frac{\det(D^{\rm ch}[q(t/e)\widehat a_0]+\bar m\eta)}
                    {\det(D^{\rm ch}[0]+\bar m\eta)},\qquad
 c_m=N_0^{(m)}(0),\quad m=2,4.
\end{equation}
The free mass is positive even when $\boldsymbol\mu_N$ is signed.
Write $\mathbb E^{(0)}_{n_*}$ for the same-mass reference law
conditioned on zero flux, and set
\begin{equation}\label{eq:vtm-Wick}
 H_2(x;v)=x^2-v,\qquad H_4(x;v)=x^4-6vx^2+3v^2,\qquad
 W_m(x;v,c)=e^{-m}H_m(x;v)-c.
\end{equation}
Here and below $e$ is the charge, so $e^{-m}$ is its inverse
power. With $X_c=\langle d\widehat a_c,F^{\rm u}\rangle$, put
\begin{equation}\label{eq:vtm-reference-root}
 \Theta_m^{(0)}
 =\left|\mathbb E_{n_*}^{(0)}
       \prod_{c\in\{0,\ldots,2L-1\}^3}W_m(X_c;v_\eta,c_m)
   \right|^{1/N_*}.
\end{equation}
Let $z_{r,n}$ be the positive zero-flux matter factor in
\eqref{eq:external-matter-factor}. At continuum mass $\bar m$ define
\[
 b(f)=-6\tr\left[
    \bigl((\bar m+\gamma\cdot\partial)^{-1}i\gamma\cdot d^*\beta
                 \bigr)^4\right],
\]
where $\gamma_\nu=\operatorname{diag}(\sigma_\nu,-\sigma_\nu)$.
The displayed fourth power is trace class and its trace is real
by \cite[Theorem~A]{PositiveProbe}.

\begin{theorem}[Volume transfer and component cumulants]\label{thm:vtm-main}
Fix the data above and a charge interval $0<e\le e_*\le1/3$.
For each charge let $N_{\rm ref}(e)$ be a finite nonnegative integer.
Assume the following conditions for every charge in this interval
and every $N\ge N_{\rm ref}(e)$.
\begin{enumerate}[label=\textup{(\roman*)}]
\item The scale and placement conditions are
\[
 v_K\ge8,\qquad e\sqrt h\log h\le1,\qquad
 6N_*Mv_K^2\exp(-v_K)
          \bigl(v_K\exp(-v_K)+\tfrac12\bigr)\le\pi^2,
\]
and
\begin{equation}\label{eq:cells-fit}
 D_\beta^{\rm buf}+3\eta\le\ell_c,\qquad
 v+Q_\beta\subset[\ell_c/8,7\ell_c/8]^3
\end{equation}
for each centered placement used.
\item There are finite nonnegative constants $C_{\rm mat}$ and
$C_{\rm ref}$, independent of charge, mesh and placement, such that,
at the same $q,\boldsymbol\mu_N$ on the two reference tori,
\begin{align}
 L^3\log z_{r,n_*}-\log z_{r,Ln_*}&\le C_{\rm mat},\notag\\
 \Theta_4^{(0)}&\le C_{\rm ref}e^{\mathfrak e},\qquad
 (1+|c_2|)\Theta_2^{(0)}\le C_{\rm ref}e^{\mathfrak e},
 \label{eq:vtm-reference-hypotheses}\\
 |c_2|+|c_4|^{1/2}&\le C_{\rm ref},\qquad
 |c_4-3c_2^2-b(f)|\le C_{\rm ref}(e+\eta^{1/4}),\notag
\end{align}
where
\[
 \mathfrak e=\tfrac12(1/4-9\varepsilon)-12\varepsilon,\qquad
 a_L=\frac{L^3-1}{(L-1)^3},\qquad
 \sigma=\mathfrak e-\frac{2a_L}{N_*}>0.
\]
\end{enumerate}
There is $C_U<\infty$, depending only on these fixed data and
constants, such that the following assertions hold uniformly
in the stated charge interval.

For each $N\ge N_{\rm ref}(e)$ there are $u_2,u_4\ge0$,
independent of the outer cutoff $T$, with
\begin{equation}\label{eq:vtm-main-errors}
 u_4\le C_Ue^\sigma,\qquad
 (1+|c_2|)u_2\le C_Ue^\sigma,\qquad
 |c_4-3c_2^2-b(f)|\le C_U(e^\sigma+\eta^{1/4}).
\end{equation}
For every $T\ge T^*(e)+1$ and every finite
$Y\subset(R^*/L)\Z^3$ whose points are incongruent modulo $R\Z^3$,
\begin{equation}\label{eq:vtm-main-products}
 \left|\E_{N,T}\prod_{y\in Y}
 W_m(F^\sharp(\tau_yf^\eta);v_\eta,c_m)\right|
              \le u_m^{|Y|},\qquad m=2,4.
\end{equation}
Here $\tau_yf(x)=f(x-y)$ and $\E_{N,T}$ is expectation under
$\mathbb P_{n;q,\boldsymbol\mu_N}^{\rm af}$. Its smooth-test
fourth cumulant satisfies
\begin{equation}\label{eq:vtm-main-finite-cumulant}
 |\kappa_4(F^\sharp(f))-e^4b(f)|
 \le(8C_U+3C_U^2)e^{4+\sigma}
                  +C_U(\eta^2+e^4\eta^{1/4}).
\end{equation}

Fix one charge. Let $\mu$ be any prescribed volume-first limit:
for $N_i\to\infty$, choose fixed-spacing weak volume limits
$\mu_i$ of these laws, with $\mu_i\Rightarrow\mu$.
All weak convergence here is on the negative Hermite space
$H_{-4}$ specified in \cref{sec:continuum-limits}.
Assume that $\mu$ is translation invariant and that almost every
component in its translation-ergodic decomposition
$\mu=\int\mu_\xi\,\dd\nu(\xi)$ satisfies
\begin{equation}\label{eq:vtm-component-G}
 \E_{\mu_\xi}\exp(F(g))\le\exp(\|g\|_2^2/2)
 \quad\text{for all real }g\in\mathcal S(\R^3;\Lambda^2).
\end{equation}
Then, on one $\nu$-full set,
\begin{equation}\label{eq:vtm-main-component}
 |\kappa_4^{\mu_\xi}(f)-e^4b(f)|
                    \le C_fe^{4+\sigma},\qquad C_f=8C_U+3C_U^2.
\end{equation}
If $b(f)>0$ and $C_fe^\sigma\le b(f)/2$, every component
in this set has $\kappa_4^{\mu_\xi}(f)\ge e^4b(f)/2>0$.
The full-measure set is for this fixed charge, probe and parent law.
\end{theorem}

For the tuned chamber construction, the same mass assignment
is selected on the top torus by \cite[Theorem~A]{CommonTuning}.
The reference-response estimate is \cite[Theorem~A]{ReferenceResponse},
and the two-torus matter comparison is \cite[Theorem~A]{TerminalPressure}.
The free coefficients and a probe with $b(f)>0$ are supplied by
\cite[Theorems~A--B]{PositiveProbe}. The volume-first laws are
constructed in \cite{ContinuumExistence}; the component Gaussian
bound is the last clause of \cite[Theorem~A]{ContinuumComponents}.
The simultaneous admission of the construction parameters is
separate from the implication proved here.

The proof is the volume-transfer and component argument of
\cite[\S25--26]{QED3Master}, with its finite-periodic moment
passage from \cite[\S6.1--6.2]{QED3Master}.
Its first difficulty is normalization. Folding bounds the
unnormalized observable by its value on the small reference torus,
but the denominator belongs to the larger torus. Sidewise
convexity of the partition function controls this denominator
using the two reference volumes. The exact photon determinant
then shows that the normalization costs only the power
$e^{-2a_L/N_*}$; the response estimate retains the positive
exponent $\sigma$.

The reference response is in flux zero. We split the full
reference expectation by flux and bound its nonzero-flux part
with the Gaussian tail and high moments. This gives bounds on
every finite product of distinct translated Wick polynomials,
uniformly in the outer volume. Their passage to a prescribed
parent uses the periodic seminorm at finite volume and the sharp
$L^2$ norm only after the supports fit in a period cube.
Finally, spatial averages and a count of repeated indices bound
the invariant conditional expectations of both Wick polynomials.
These two conditional bounds determine the fourth cumulant in
each component.

The charge restriction for positivity is imposed after $C_U$:
\begin{equation}\label{eq:constants-probe-charge}
 0<e_\beta\le
       \min\{e_*,[b(f)/(2C_f)]^{1/\sigma}\}.
\end{equation}
It is positive when $b(f)>0$ and $C_f>0$.
The conclusions concern subsequential limits of the stated
cutoff family; no uniqueness of the continuum law is required.

